% GENERATED FILE. Edit canonical lessons, then run npm run generate:books.
% canonical-source-hash: fnv1a64:f40723f91c72a55f
% canonical-lessons: AR-C546-shawt, AR-C546-mushajji, AR-C546-taadala, AR-C546-marrara, AR-C546-tanafasa

\chapter{Match day}
\label{ch:ar-match-day}
\hlchaptermodality{546}

\begin{chapteropening}
\textbf{By the end of this chapter:} \emph{I can talk about a half and a supporter, and say that teams finished level, competed or passed the ball.}

Complete the last lesson of chapter 546: I can talk about a half and a supporter, and say that teams finished level, competed or passed the ball.
\end{chapteropening}

\section[\texorpdfstring{\ar{شوط} (shawṭ)}{شوط (shawṭ)}]{\ar{شوط} (shawṭ) — a half (of a match)}

\label{lesson:AR-C546-shawt}



\begin{quote}
Before the new one: say the Arabic for to create, then the Arabic for to direct films.
\end{quote}

\subsection*{You'll want to know: \ar{شوط}}
\textbf{\ar{شوط}} — \emph{shawṭ} — \textquotedblleft{}a half (of a match)\textquotedblright{}.

\begin{quote}
\ar{سجل} \ar{الفريق} \ar{هدفا} \ar{في} \ar{الشوط} \ar{الأول}.
\end{quote}

\begin{quote}
\emph{sajjala al-farīq hadafan fī ash-shawṭ al-awwal.}
\end{quote}

\begin{quote}
The team scored a goal in the first half.
\end{quote}

Plural: \ar{أشواط}, \emph{ashwāṭ}.

\begin{cousinweb}
\textbf{\ar{شوط}} (\emph{shawṭ}) is built on the native Arabic root \emph{sh-w-ṭ}, \textquotedblleft{}a run, a single course\textquotedblright{} (Lane). Lane gives shawṭ as one run of a course: one stretch of a race.
\end{cousinweb}

\begin{grammarlens}[title={What you've built}]
One more word for describing what you lived and felt.
\end{grammarlens}

\subsection*{Your turn}
\begin{itemize}
\raggedright
  \item \emph{Say it:} \emph{shawṭ}
  \item \emph{Say it:} \emph{shawṭ}, once more
  \item \emph{Say it:} \emph{sajjala al-farīq hadafan fī ash-shawṭ al-awwal.} — The team scored a goal in the first half.
  \item \emph{Recall:} say the Arabic for to create, then the Arabic for to direct films, then say \emph{shawṭ} again
\end{itemize}

\begin{tcolorbox}[breakable,colback=teal!4,colframe=teal!35!black,title={Before you move on}]
What does \ar{شوط} mean? (\textquotedblleft{}A half (of a match)\textquotedblright{}.) Say it once more.
\end{tcolorbox}
\section[\texorpdfstring{\ar{مشجع} (mushajjiʿ)}{مشجع (mushajjiʿ)}]{\ar{مشجع} (mushajjiʿ) — a supporter, a fan}

\label{lesson:AR-C546-mushajji}



\begin{quote}
Before the new one: say the Arabic for to direct films, then the Arabic for a half.
\end{quote}

\subsection*{You'll want to know: \ar{مشجع}}
\textbf{\ar{مشجع}} — \emph{mushajjiʿ} — \textquotedblleft{}a supporter, a fan\textquotedblright{}.

\begin{quote}
\ar{غنى} \ar{المشجعون} \ar{في} \ar{الملعب}.
\end{quote}

\begin{quote}
\emph{ghannā al-mushajjiʿūn fī al-malʿab.}
\end{quote}

\begin{quote}
The supporters sang in the stadium.
\end{quote}

\begin{cousinweb}
\textbf{\ar{مشجع}} (\emph{mushajjiʿ}) is built on the native Arabic root \emph{sh-j-ʿ}, \textquotedblleft{}to be brave\textquotedblright{} (Lane). Mushajjiʿ is the participle of shajjaʿa, to encourage.
\end{cousinweb}

\begin{grammarlens}[title={What you've built}]
One more word for describing what you lived and felt.
\end{grammarlens}

\subsection*{Your turn}
\begin{itemize}
\raggedright
  \item \emph{Say it:} \emph{mushajjiʿ}
  \item \emph{Say it:} \emph{mushajjiʿ}, once more
  \item \emph{Say it:} \emph{ghannā al-mushajjiʿūn fī al-malʿab.} — The supporters sang in the stadium.
  \item \emph{Recall:} say the Arabic for to direct films, then the Arabic for a half, then say \emph{mushajjiʿ} again
\end{itemize}

\begin{tcolorbox}[breakable,colback=teal!4,colframe=teal!35!black,title={Before you move on}]
What does \ar{مشجع} mean? (\textquotedblleft{}A supporter, a fan\textquotedblright{}.) Say it once more.
\end{tcolorbox}
\section[\texorpdfstring{\ar{تعادل} (taʿādala)}{تعادل (taʿādala)}]{\ar{تعادل} (taʿādala) — to finish level; to end in a draw}

\label{lesson:AR-C546-taadala}



\begin{quote}
Before the new one: say the Arabic for a half, then the Arabic for a supporter.
\end{quote}

\subsection*{You'll want to know: \ar{تعادل}}
\textbf{\ar{تعادل}} — \emph{taʿādala} — \textquotedblleft{}to finish level; to end in a draw\textquotedblright{}.

\begin{quote}
\ar{تعادل} \ar{الفريقان} \ar{في} \ar{المباراة}.
\end{quote}

\begin{quote}
\emph{taʿādala al-farīqān fī al-mubārā.}
\end{quote}

\begin{quote}
The two teams drew the match.
\end{quote}

Present: \ar{يتعادل}, \emph{yataʿādal}.

\begin{cousinweb}
\textbf{\ar{تعادل}} (\emph{taʿādala}) is built on the native Arabic root \emph{ʿ-d-l}, \textquotedblleft{}to be equal, to be fair\textquotedblright{} (Lane). Form VI taʿādala is to be equal to one another.
\end{cousinweb}

\begin{grammarlens}[title={What you've built}]
One more word for describing what you lived and felt.
\end{grammarlens}

\subsection*{Your turn}
\begin{itemize}
\raggedright
  \item \emph{Say it:} \emph{taʿādala}
  \item \emph{Say it:} \emph{taʿādala}, once more
  \item \emph{Say it:} \emph{taʿādala al-farīqān fī al-mubārā.} — The two teams drew the match.
  \item \emph{Recall:} say the Arabic for a half, then the Arabic for a supporter, then say \emph{taʿādala} again
\end{itemize}

\begin{tcolorbox}[breakable,colback=teal!4,colframe=teal!35!black,title={Before you move on}]
What does \ar{تعادل} mean? (\textquotedblleft{}To finish level; to end in a draw\textquotedblright{}.) Say it once more.
\end{tcolorbox}
\section[\texorpdfstring{\ar{مرر} (marrara)}{مرر (marrara)}]{\ar{مرر} (marrara) — to pass the ball}

\label{lesson:AR-C546-marrara}



\begin{quote}
Before the new one: say the Arabic for a supporter, then the Arabic for to finish level.
\end{quote}

\subsection*{You'll want to know: \ar{مرر}}
\textbf{\ar{مرر}} — \emph{marrara} — \textquotedblleft{}to pass the ball\textquotedblright{}.

\begin{quote}
\ar{مرر} \ar{اللاعب} \ar{الكرة} \ar{إلى} \ar{زميله}.
\end{quote}

\begin{quote}
\emph{marrara al-lāʿib al-kura ilā zamīlihi.}
\end{quote}

\begin{quote}
The player passed the ball to his teammate.
\end{quote}

Present: \ar{يمرر}, \emph{yumarrir}.

\begin{cousinweb}
\textbf{\ar{مرر}} (\emph{marrara}) is built on the native Arabic root \emph{m-r-r}, \textquotedblleft{}to pass, to go by\textquotedblright{} (Lane). Form II marrara is to make a thing pass.
\end{cousinweb}

\begin{grammarlens}[title={What you've built}]
One more word for describing what you lived and felt.
\end{grammarlens}

\subsection*{Your turn}
\begin{itemize}
\raggedright
  \item \emph{Say it:} \emph{marrara}
  \item \emph{Say it:} \emph{marrara}, once more
  \item \emph{Say it:} \emph{marrara al-lāʿib al-kura ilā zamīlihi.} — The player passed the ball to his teammate.
  \item \emph{Recall:} say the Arabic for a supporter, then the Arabic for to finish level, then say \emph{marrara} again
\end{itemize}

\begin{tcolorbox}[breakable,colback=teal!4,colframe=teal!35!black,title={Before you move on}]
What does \ar{مرر} mean? (\textquotedblleft{}To pass the ball\textquotedblright{}.) Say it once more.
\end{tcolorbox}
\section[\texorpdfstring{\ar{تنافس} (tanāfasa)}{تنافس (tanāfasa)}]{\ar{تنافس} (tanāfasa) — to compete}

\label{lesson:AR-C546-tanafasa}



\begin{quote}
Before the new one: say the Arabic for to finish level, then the Arabic for to pass the ball.
\end{quote}

\subsection*{You'll want to know: \ar{تنافس}}
\textbf{\ar{تنافس}} — \emph{tanāfasa} — \textquotedblleft{}to compete\textquotedblright{}.

\begin{quote}
\ar{تنافس} \ar{الفريقان} \ar{على} \ar{البطولة}.
\end{quote}

\begin{quote}
\emph{tanāfasa al-farīqān ʿalā al-buṭūla.}
\end{quote}

\begin{quote}
The two teams competed for the championship.
\end{quote}

Present: \ar{يتنافس}, \emph{yatanāfas}.

It takes \ar{على} for the prize.

\begin{cousinweb}
\textbf{\ar{تنافس}} (\emph{tanāfasa}) is built on the native Arabic root \emph{n-f-s}, \textquotedblleft{}a soul; to be precious\textquotedblright{} (Lane). Form VI tanāfasa is to vie with one another for a precious thing.
\end{cousinweb}

\begin{grammarlens}[title={What you've built}]
One more word for describing what you lived and felt.
\end{grammarlens}

\subsection*{Your turn}
\begin{itemize}
\raggedright
  \item \emph{Say it:} \emph{tanāfasa}
  \item \emph{Say it:} \emph{tanāfasa}, once more
  \item \emph{Say it:} \emph{tanāfasa al-farīqān ʿalā al-buṭūla.} — The two teams competed for the championship.
  \item \emph{Recall:} say the Arabic for to finish level, then the Arabic for to pass the ball, then say \emph{tanāfasa} again
\end{itemize}

\begin{tcolorbox}[breakable,colback=teal!4,colframe=teal!35!black,title={Before you move on}]
What does \ar{تنافس} mean? (\textquotedblleft{}To compete\textquotedblright{}.) Say it once more.
\end{tcolorbox}
